\chapter{Tail Hedging and Long Volatility}\label{ch:s2:tail-hedging-and-long-volatility}

A \omterm{def:s2:tail-hedging-and-long-volatility:hedge}{tail hedge} that bleeds a little every month is the most disliked position in a portfolio for years, until the month it pays for all of them. Cboe's
5\% Put Protection Index, which holds the S\&P 500 and buys a 5\% out-of-the-money put every month, lost 10.6\% in October 1987 when the index fell
21.8\%, and 0.4\% over February and March 2020 when the index fell 19.9\%. From July 2007 to March 2009 it still lost 42.0\%, because monthly puts
protect against a crash but not against a slow bear market. On this chapter's synthetic market a monthly put programme costs 6.0\% a year in years
without a crash, pays 19.2\% and 39.5\% in two of the three crash years, and leaves the portfolio growing at 2.5\% a year against 3.0\% unhedged,
with its worst drawdown cut from 73\% to 41\%. Whether that was worth it depends on what the hedge lets the rest of the portfolio do. The build is
\texttt{firm.tailhedge}.

\section{The cost of convexity}

\begin{definition}[Tail hedge]\label{def:s2:tail-hedging-and-long-volatility:hedge}
A \emph{tail hedge}\index{tail hedge} is a position held to pay off in large, rare market falls, such as out-of-the-money puts, put spreads, long
variance or volatility futures, bought at a cost in ordinary times in exchange for convexity in a crash.
\end{definition}

\begin{definition}[Hedge bleed]\label{def:s2:tail-hedging-and-long-volatility:bleed}
The \emph{hedge bleed}\index{hedge bleed} is a \omterm{def:s2:tail-hedging-and-long-volatility:hedge}{tail hedge}'s cost in periods without a crash: premiums paid less payoffs received, as a share of the
portfolio, per year.
\end{definition}

An out-of-the-money put is priced at its point on the smile, which for index puts is above the at-the-money vol. That vol already carries the variance
premium of chapter~1. The buyer pays both premiums, the one the seller of chapter~1 collects. Harvey and co-authors compared defensive strategies over
large equity drawdowns and recessions. Continuously holding short-dated S\&P 500 puts was the most reliable defence and also the most costly.

\texttt{firm.tailhedge} runs put programmes on the synthetic index of chapter~1. It is twenty years long, with three planted crashes of 20\% to 30\%
in ten days, in years 5, 12 and 17. Options are priced on a steeper smile than chapter~1's (skew $-0.25$ per standardised unit, curvature 0.02). The
portfolio holds the index and one put per unit of index, and at each roll it sets its units so that equity and puts cover the same notional
(\cref{lst:s2:tail-hedging-and-long-volatility:programme}).

\section{Put programmes and their monetisation}

\begin{definition}[Hedge monetisation]\label{def:s2:tail-hedging-and-long-volatility:monetisation}
\emph{Hedge monetisation}\index{hedge monetisation} is the rule for selling a \omterm{def:s2:tail-hedging-and-long-volatility:hedge}{tail hedge} that has gained, before its expiry, and replacing it at the
new market level: for instance, when the put is worth three times what it cost. It turns a paper gain into cash that can be reinvested, and
restores protection near the new spot.
\end{definition}

Cboe's indices show the real shape (\cref{fig:s2:tail-hedging-and-long-volatility:episodes}). From June 1986 to September 2026, PPUT, which is a
total-return index, grew at a log rate of 7.7\% a year with a volatility of 13.5\%. CLLZ holds the S\&P 500, buys a 2.5\%--5\% put spread each month
and pays for it by selling a call. It grew at 7.8\% with a volatility of 14.6\%. The S\&P 500 price index grew at 8.5\% with a volatility of 18.4\%.
The price index omits dividends, so the return gap understates the hedges' cost. PPUT's worst 21 days cost 13.5\%, the index's 33.0\%. In the
crash of autumn 2008, September to November, PPUT lost 6.7\%, the index 30.1\% and CLLZ 31.2\%: a put spread protects only the fall between its
strikes. The largest drawdowns tell the rest. PPUT lost 42.0\% from July 2007 to March 2009, when the index lost 56.8\%. The fall took twenty
months, and a put struck 5\% below each month's start pays only for the part of that month's fall beyond 5\%.

\begin{omfigure}
\begin{tikzpicture}
\begin{axis}[width=11cm,height=5.6cm,ybar,bar width=8pt,ylabel={\small return (\%)},tick label style={font=\footnotesize},
  xmin=0.5,xmax=4.5,ymin=-60,ymax=5,xtick={1,2,3,4},xticklabels={Oct 1987,autumn 2008,early 2020,max drawdown},
  grid=major,grid style={gray!25},table/col sep=comma,legend style={font=\scriptsize,at={(0.5,-0.2)},anchor=north},legend columns=3]
\addplot[fill=omThm!70,draw=omThm] table[x=x,y=PPUT]{figdata/strategies-2/07-tail-hedging-and-long-volatility/episodes.csv};
\addplot[fill=omDef!60,draw=omDef] table[x=x,y=CLLZ]{figdata/strategies-2/07-tail-hedging-and-long-volatility/episodes.csv};
\addplot[fill=gray!40,draw=gray] table[x=x,y=SPX]{figdata/strategies-2/07-tail-hedging-and-long-volatility/episodes.csv};
\legend{PPUT (5\% monthly puts),CLLZ (put-spread collar),S\&P 500 price index}
\end{axis}
\end{tikzpicture}

\omcaption{Cboe's put-protection and collar indices against the S\&P 500 in three crashes and over each series' largest drawdown (all three from
2007 to March 2009), 1986--2026. Monthly puts protect against fast falls, not slow ones. Derived from Cboe's index histories; the raw series are not
redistributed. Data: \texttt{s2\_fetch\_protect}.}\label{fig:s2:tail-hedging-and-long-volatility:episodes}
\end{omfigure}

On the synthetic market the programmes compare as follows, each against the index and against the simplest alternative, a static 70\% in the index
and 30\% in cash rebalanced monthly:

\begin{center}
\small
\begin{tabular}{@{}lcccc@{}}
\toprule
twenty years & growth (\%/yr) & volatility (\%) & max drawdown (\%) & worst month (\%)\\
\midrule
index & 3.0 & 19.0 & $-73$ & $-43$\\
5\% puts, monthly & 2.5 & 13.9 & $-41$ & $-17$\\
10\% puts, monthly & 3.2 & 16.0 & $-49$ & $-22$\\
10\% puts, quarterly & 3.0 & 14.0 & $-46$ & $-24$\\
20\% puts, quarterly & 4.2 & 16.2 & $-52$ & $-33$\\
10\% quarterly, monetised at 3$\times$ & 2.2 & 14.0 & $-62$ & $-26$\\
20\% puts, half-yearly & 4.6 & 15.5 & $-50$ & $-36$\\
70\% index, 30\% cash & 2.8 & 13.2 & $-58$ & $-31$\\
index + trend overlay & 6.3 & 17.6 & $-56$ & $-24$\\
\bottomrule
\end{tabular}
\end{center}

The yearly accounting is the hedge's own. Over the 17 calm years, the monthly 5\% programme bled 6.0\% of the portfolio a year. In the three crash
years it earned 19.2\%, lost 2.7\% and earned 39.5\%. The second crash came when the index's vol stood at 3.9\% and the put was cheap, and the put
still paid only 12.0\% of the portfolio in the crash's own weeks, against 30.0\% and 29.2\% in the other two. After a crash vol stays high for months, and the puts bought then are dear: that year's bleed consumed the
payoff. Farther puts bleed less (1.5\% a year for 20\% half-yearly puts) and pay later and less (2.0\%, $-2.1\%$ and 26.3\%). The monetised programme
sold its puts eleven times at three times their cost. It earned the most in the first crash year (16.3\%), selling early in the fall. It then held only protection struck near
the new, lower spot, and the rest of the fall passed below it: its drawdown, 62\%, is the worst of the programmes.

\section{Alternatives: trend, long vol, variance}

Harvey and co-authors found that futures time-series momentum did well in extended sell-offs and a quality long--short strategy in flights to
quality. Treasuries carry positively, but the post-2000 negative correlation between bonds and equities is a historical rarity. Long gold and long
credit protection sat between puts and bonds in cost and reliability. On the synthetic index, a trend overlay adds half the index's notional in the
direction of its trailing year's return. It grows at 6.3\% a year with a 56\% drawdown and a worst month of $-24\%$: better than the puts on growth,
worse on drawdown. It needs a sell-off long enough for the signal to turn. Crisis alpha of that kind (Book~8, chapter~19) and the convexity of an
option insure different things.

Long variance and long VIX futures are the volatility book's versions of the hedge. They pay the variance premium and the roll-down of chapter~5 in
ordinary times, and pay out in a spike of volatility whether or not the index gaps. They need their own monetisation rule, since a variance position's
gain can disappear within weeks once volatility falls back.

\section{Judging a hedge by the portfolio}

Israelov's warning is the test to apply. In the typical use, puts were quite ineffective at reducing drawdowns compared with the simple alternative
of statically holding less of the index. Unless purchases and maturities are timed right around the drawdown, they may add to drawdowns and to
volatility per unit of expected return. The synthetic market is kind to puts: its crashes are fast, ten days each, which is the case a put is built
for. There the 5\% monthly programme beats the static mix on drawdown ($-41\%$ against $-58\%$ at similar volatility) and loses on growth (2.5\% against
2.8\%). The real index's worst drawdown was slow, and PPUT lost 42.0\% of it.

So a hedge is judged by what it lets the portfolio do. With a hedge that caps the loss, the rest of the portfolio can hold more risk, so the cost
comes back as extra exposure. Without that, the hedge is a drag with a good story. The growth rate is the measure, over a history long enough to
contain crashes of both kinds.

\begin{omfigure}
\begin{tikzpicture}
\begin{axis}[width=11cm,height=5.6cm,ymode=log,xlabel={\small year},ylabel={\small value (start 1)},tick label style={font=\footnotesize},
  xmin=0,xmax=20,ymin=0.7,ymax=5,ytick={0.75,1,1.5,2,3,4},log ticks with fixed point,grid=major,grid style={gray!25},table/col sep=comma,
  legend style={font=\scriptsize,at={(0.5,-0.3)},anchor=north},legend columns=4]
\addplot[gray,thick] table[x=year,y=index]{figdata/strategies-2/07-tail-hedging-and-long-volatility/paths.csv};
\addplot[omThm,thick] table[x=year,y=puts]{figdata/strategies-2/07-tail-hedging-and-long-volatility/paths.csv};
\addplot[omDef,thick,dashed] table[x=year,y=static]{figdata/strategies-2/07-tail-hedging-and-long-volatility/paths.csv};
\addplot[black,thin] table[x=year,y=trend]{figdata/strategies-2/07-tail-hedging-and-long-volatility/paths.csv};
\legend{index,5\% monthly puts,70\% index + cash,index + trend}
\end{axis}
\end{tikzpicture}

\omcaption{The synthetic index and three defended portfolios over twenty years with crashes in years 5, 12 and 17, on a log scale: a monthly
5\% put programme, a static 70\% in the index, and the index with a trend overlay. Data: \texttt{s2\_tailhedge.curves}.}\label{fig:s2:tail-hedging-and-long-volatility:paths}
\end{omfigure}

\section{Strategy files}

\begin{strategyfile}{Rolling out-of-the-money put programme}\label{strat:s2:tail-hedging-and-long-volatility:puts}
\sfield{Who pays you, and why} Nobody, on average: the buyer pays the variance and skew premiums for convexity in a crash.
\sfield{Instruments and venues} Listed index puts, monthly to half-yearly, 5\% to 20\% out of the money.
\sfield{Signal} None; a programme, not a trade. Some programmes buy more when protection is cheap.
\sfield{Sizing and execution} A fixed budget a year, rolled mechanically; monetisation rules decided in advance.
\sfield{Costs} Premiums and spreads; the bleed.
\sfield{How it dies} Slow bear markets that fall within each put's strike; being abandoned after years of bleed, just before it pays.
\sfield{Horizon, capacity, infrastructure} Years; large capacity in index options.
\sfield{Backtest honestly} Include slow and fast drawdowns; compare with statically holding less.
\sfield{Sources} Israelov (2019); Harvey and co-authors (2019); Cboe PPUT: $-42.0\%$ from July 2007 to March 2009; this chapter: a 6.0\% bleed a year
and a drawdown cut from 73\% to 41\%.
\end{strategyfile}

\begin{strategyfile}{Put spread collar}\label{strat:s2:tail-hedging-and-long-volatility:collar}
\sfield{Who pays you, and why} The call buyer finances the put spread; the holder gives up upside to cap a band of downside.
\sfield{Instruments and venues} Index put spreads and calls, monthly.
\sfield{Signal} None; a programme.
\sfield{Sizing and execution} Zero-cost by construction: the call's strike set so that its premium pays for the put spread.
\sfield{Costs} Spreads on three options a month; the upside given away.
\sfield{How it dies} A crash beyond the put spread's lower strike: CLLZ lost 31.2\% in autumn 2008.
\sfield{Horizon, capacity, infrastructure} Months; a mechanical roll.
\sfield{Backtest honestly} Measure crashes below the lower strike, not only the band.
\sfield{Sources} Cboe CLLZ (derived statistics in this chapter).
\end{strategyfile}

\begin{strategyfile}{Long variance with monetisation}\label{strat:s2:tail-hedging-and-long-volatility:variance}
\sfield{Who pays you, and why} Nobody, on average: it pays the variance premium for a payoff in volatility spikes.
\sfield{Instruments and venues} Variance swaps, VIX futures or calls.
\sfield{Signal} Cheap volatility against its history; the slope of the VIX curve.
\sfield{Sizing and execution} A budget; monetise at a set multiple or when volatility has spiked and the curve inverts.
\sfield{Costs} The variance premium; the roll-down of VIX futures.
\sfield{How it dies} Spikes that fade before monetisation; a crash with little volatility.
\sfield{Horizon, capacity, infrastructure} Months; variance or VIX access.
\sfield{Backtest honestly} Use the price at the monetisation time, not the peak.
\sfield{Sources} Harvey and co-authors (2019) on defensive strategies; no performance figure verified.
\end{strategyfile}

\section{Tutorial: paying for the month that pays}\label{tut:s2:tail-hedging-and-long-volatility:tut}

\begin{tutorial}
\textbf{Goal.} Run put programmes of several strikes, tenors and monetisation rules on the synthetic index, compare them with a static mix and a trend
overlay, and read Cboe's protection indices. \textbf{End state:} the table and the two figures.
\begin{tutsteps}
\item \textbf{The programme}.
\omcode{code/firm/tailhedge/firm_tailhedge.py}{43}{63}{A rolling put programme with monetisation, fully invested.}\label{lst:s2:tail-hedging-and-long-volatility:programme}
\item \textbf{The bleed}.
\omcode{code/strategies-2/07-tail-hedging-and-long-volatility/python/s2_tailhedge.py}{57}{72}{Each programme's yearly hedge P\&L, calm years and crash years.}\label{lst:s2:tail-hedging-and-long-volatility:bleed}
\item \textbf{Run} \texttt{s2\_fetch\_protect.py} once, then \texttt{table()}, \texttt{bleed()} and \texttt{fig\_tailhedge.py}.
\end{tutsteps}
\textbf{What to change next.} Plant a slow bear market (a year of steady decline) and rerun; lever the hedged portfolio to the unhedged one's
volatility and compare growth; buy puts only when the skew is flat.
\end{tutorial}

\section{Build: tail hedge}\label{bld:s2:tail-hedging-and-long-volatility:tailhedge}

\begin{build}
\bfield{Purpose} Rolling put programmes with monetisation, static mixes and trend overlays, judged by growth, volatility and drawdown.
\bfield{Interface} \texttt{put\_vol(S, K, tau, v, cfg)}, \texttt{put\_programme(r, v, cfg, otm, tenor, monetise)}, \texttt{static\_mix(r, weight,
every)}, \texttt{trend\_overlay(r, weight, lookback)}, \texttt{evaluate(nav)}.
\bfield{Rules} Puts priced and marked on the surface with the smile clipped; fully invested at each roll; monetisation replaces the put at the new
spot.
\bfield{Acceptance tests} \texttt{code/firm/tailhedge/tests/}: a still market only bleeds; the put covers a fall below its strike; a static mix, a
trend overlay and a drawdown by hand.
\bfield{Stretch} Put spreads and collars; variance and VIX hedges; leverage to a target volatility.
\end{build}

\begin{omsources}
\item R.~Israelov, ``Pathetic protection: the elusive benefits of protective puts'', \emph{Journal of Alternative Investments} 21(3), 2019.
\item C.~R.~Harvey, E.~Hoyle, S.~Rattray, M.~Sargaison, D.~Taylor and O.~Van Hemert, ``The best of strategies for the worst of times: can portfolios be
crisis proofed?'', \emph{Journal of Portfolio Management} 45(5), 2019.
\item Cboe, S\&P 500 Put Protection Indices methodology; Cboe Insights on PPUT, CLL and CLLZ, 2021.
\item Cboe Global Markets, daily histories of PPUT, CLLZ and SPX.
\end{omsources}

\section{Exercises}

\begin{exercise}[$\star$]\label{exo:s2:tail-hedging-and-long-volatility:1}
A put programme bleeds 6.0\% a year in calm years. Over 17 calm years and three crash years paying 19.2\%, $-2.7\%$ and 39.5\%, what is its total,
ignoring compounding?
\end{exercise}

\begin{exercise}[$\star$]\label{exo:s2:tail-hedging-and-long-volatility:2}
Why does a monthly 5\% put fail to protect against a fall of 1\% a week for a year?
\end{exercise}

\begin{exercise}[$\star$]\label{exo:s2:tail-hedging-and-long-volatility:3}
A put costs 1\% of notional and is monetised at three times its cost. What is it worth when sold?
\end{exercise}

\begin{exercise}[$\star\star$]\label{exo:s2:tail-hedging-and-long-volatility:4}
Why did CLLZ lose almost as much as the index in autumn 2008 while PPUT lost 6.7\%?
\end{exercise}

\begin{exercise}[$\star\star$]\label{exo:s2:tail-hedging-and-long-volatility:5}
Why is the return gap between PPUT and the S\&P 500 price index an understatement of the hedge's cost?
\end{exercise}

\begin{exercise}[$\star\star$]\label{exo:s2:tail-hedging-and-long-volatility:6}
How does a trend overlay's protection differ from a put's?
\end{exercise}

\begin{exercise}[$\star\star\star$]\label{exo:s2:tail-hedging-and-long-volatility:7}
\emph{Coding.} Run \texttt{slow\_bear()}, which adds a steady decline of 40\% in log terms over year 8 of the synthetic index, and compare the
year's loss of the 5\% monthly programme, the static mix and the index. Which protects better, and why?
\end{exercise}

\begin{exercise}[$\star\star\star$]\label{exo:s2:tail-hedging-and-long-volatility:8}
\emph{Find the flaw.} ``The hedge cost us 6\% a year for five years; it does not work, so we are removing it.''
\end{exercise}

\section{Problem: Paying for the Month That Pays}

\begin{problem}[{Weekend problem --- a tail-hedging programme}]\label{pb:s2:tail-hedging-and-long-volatility:1}
The chapter's synthetic market, Cboe's protection indices and the public record.

\medskip\noindent\textbf{Part I --- The cost.}
\begin{enumerate}
\item Define a \omterm{def:s2:tail-hedging-and-long-volatility:hedge}{tail hedge} and its bleed.
\item Which premiums does a put buyer pay?
\item What did Harvey and co-authors find about puts and other defences?
\item What did Israelov find?
\end{enumerate}

\medskip\noindent\textbf{Part II --- The real indices.}
\begin{enumerate}[resume]
\item Give PPUT's, CLLZ's and the S\&P 500's growth and volatility.
\item How did each fare in 1987, 2008 and 2020?
\item Why did PPUT lose 42\% from 2007 to 2009?
\item Why does a put spread fail in a deep crash?
\end{enumerate}

\medskip\noindent\textbf{Part III --- The programmes.}
\begin{enumerate}[resume]
\item Define monetisation.
\item Compare the programmes' growth and drawdowns.
\item Give the 5\% monthly programme's bleed and crash-year payoffs.
\item Why did monetising worsen the drawdown?
\end{enumerate}

\medskip\noindent\textbf{Part IV --- The verdict.}
\begin{enumerate}[resume]
\item State the \emph{named result}: the hedge's annual bleed and the portfolio growth rate with and without it.
\item How does the hedge compare with a static mix?
\item How does it compare with a trend overlay?
\item Why is the synthetic market kind to puts?
\item How should a hedge be judged?
\item How would you backtest a \omterm{def:s2:tail-hedging-and-long-volatility:hedge}{tail hedge} honestly?
\item Which strategy file suits a slow bear market best?
\item In one sentence: what does a \omterm{def:s2:tail-hedging-and-long-volatility:hedge}{tail hedge} buy?
\end{enumerate}
\end{problem}

\section{Interview questions}

\begin{interviewq}[$\star$ \iqroles{trader}]\label{iq:s2:tail-hedging-and-long-volatility:1}
Why are out-of-the-money index puts expensive?
\end{interviewq}

\begin{interviewq}[$\star\star$ \iqroles{researcher}]\label{iq:s2:tail-hedging-and-long-volatility:2}
How would you compare a put programme with holding less equity?
\end{interviewq}

\begin{interviewq}[$\star\star$ \iqroles{trader}]\label{iq:s2:tail-hedging-and-long-volatility:3}
Your puts are worth five times what you paid. Do you sell them?
\end{interviewq}

\begin{interviewq}[$\star\star$ \iqroles{risk}]\label{iq:s2:tail-hedging-and-long-volatility:4}
What scenarios would you use to test a tail-hedging programme?
\end{interviewq}

\begin{interviewq}[$\star\star$ \iqroles{developer}]\label{iq:s2:tail-hedging-and-long-volatility:5}
Design the daily process for a rolling put programme with monetisation rules.
\end{interviewq}

\begin{interviewq}[$\star\star\star$ \iqroles{researcher}]\label{iq:s2:tail-hedging-and-long-volatility:6}
Show that for a portfolio with log-normal returns the growth rate is approximately $\mu - \sigma^2/2$, and explain how a hedge that lowers $\sigma$ can
raise the growth rate even if it lowers $\mu$.
\end{interviewq}
